% !TEX program = xelatex
% 中文 Beamer 示例 —— 使用 ctexbeamer + Madrid 主题 + beaver 配色
% 编译：xelatex beamer_demo.tex （运行两次以生成目录）
\documentclass[aspectratio=169]{ctexbeamer}
\usetheme{Madrid}
\usecolortheme{beaver}

\title{Beamer 快速上手}
\subtitle{一个简短的中文演示示例}
\author{Yuki}
\institute{小女仆工作室}
\date{\today}

\begin{document}

% ---------- 1. 标题页 ----------
\begin{frame}
  \titlepage
\end{frame}

% ---------- 2. 带 itemize 的大纲页 ----------
\begin{frame}{大纲}
  \tableofcontents
  \vspace{0.5em}
  \begin{itemize}
    \item 数学公式展示
    \item 分块（block）的使用
    \item 两栏（columns）排版
  \end{itemize}
\end{frame}

\section{数学公式}

% ---------- 3. 带公式的 frame ----------
\begin{frame}{经典数学公式}
  \begin{itemize}
    \item 质能方程：
      \begin{equation}
        E = mc^2
      \end{equation}
    \item 勾股定理：
      \begin{equation}
        a^2 + b^2 = c^2
      \end{equation}
  \end{itemize}
\end{frame}

\section{分块与两栏}

% ---------- 4. 带 block 的 frame ----------
\begin{frame}{分块（Block）示例}
  \begin{block}{定义}
    若 $a^2 + b^2 = c^2$，则以 $a$、$b$ 为直角边、$c$ 为斜边的三角形是直角三角形。
  \end{block}
  \begin{alertblock}{注意}
    公式必须放在数学模式中，例如 $E = mc^2$。
  \end{alertblock}
  \begin{exampleblock}{例子}
    取 $a = 3$、$b = 4$，则 $c = \sqrt{3^2 + 4^2} = 5$。
  \end{exampleblock}
\end{frame}

% ---------- 5. 带 columns 的 frame ----------
\begin{frame}{两栏（Columns）排版}
  \begin{columns}[T]
    \begin{column}{0.48\textwidth}
      \begin{block}{左栏：文字}
        Beamer 可以轻松实现左右分栏。
        \begin{itemize}
          \item 左栏适合放要点
          \item 结构清晰、重点突出
        \end{itemize}
      \end{block}
    \end{column}
    \begin{column}{0.48\textwidth}
      \begin{block}{右栏：公式}
        质能方程：\[ E = mc^2 \]
        勾股定理：\[ a^2 + b^2 = c^2 \]
      \end{block}
    \end{column}
  \end{columns}
\end{frame}

% ---------- 6. 结束页 ----------
\begin{frame}
  \centering
  {\Huge 谢谢观看！}
  \vspace{1em}
  {\Large Questions?}
\end{frame}

\end{document}
